\section{Visualisations}
\label{sec:charts}

Eleven of the 45 points are for the charts. Each chart below answers one question a user has about
the model's output; none is decoration.

\subsection{The charts}
\label{sec:chartlist}

{\small
\begin{longtable}{@{}L{0.7cm}L{2.7cm}L{2.0cm}L{4.6cm}L{1.5cm}L{1.6cm}@{}}
\toprule
\textbf{No.} & \textbf{Chart} & \textbf{Type} & \textbf{Question it answers} & \textbf{Data} & \textbf{Tier} \\
\midrule
\endhead
G1 & Daily activity & Bar, on a date axis & When was this address active, and on which days did something unusual happen? & E5 & \must \\
G2 & Transactions over time & Scatter, log amount & How large was each transaction, and where do the unusual ones sit? & E6 & \must \\
G3 & Feature deviation & Horizontal bar & How far is this transaction from the address's usual size, pace, gap and ratio? & E7, E8 & \must \\
G5 & What-if response & Line, log amount & How does the score change with the amount, and where does it cross the flag line? & E9 & \must \\
G4 & Model verdicts & Dot plot & Where does this transaction rank among the address's own transactions for each model, and which models flag it? & E7, E8 & \should{} S-02 \\
G6 & Planted-test recall & Grouped bar with intervals & How often does each method catch each planted pattern? & E11 & \should{} S-03 \\
G7 & TxGPT surprisal by field & Stacked bar or heat map & Which part of the transaction surprised the language model? & E21 & \could{} C-01 \\
G8 & Activity calendar & Heat map & Which weeks and weekdays is the address active? & E5 & \could{} C-03 \\
G9 & Snapshot map & Scatter, WebGL & Where are the unusual transactions across all 7,966 in the snapshot? & E6 & \could{} C-03 \\
\bottomrule
\end{longtable}
}

The Must tier alone has three chart types (bar, scatter, line) in four charts. If one Should chart
slips, the count of types does not fall.

\subsection{The chart contract}
\label{sec:chartcontract}

Every chart is drawn by one wrapper component, \fp{src/charts/PlotlyChart.tsx}, so these hold for
all of them and are tested once.

\begin{itemize}
  \item \textbf{Words on the chart (NFR-01, condition 1).} A title; both axes labelled with units
    (BTC, transactions, date in UTC, robust z); a legend when there is more than one series; and a
    caption with the source and the snapshot date, or the fetch time for live data.
  \item \textbf{Tooltips.} The facts in plain words and real units: ``28 Sep 2026, 06:51 UTC;
    received 0.00056497 BTC; unusual: longer gap than usual''. Amounts to eight places; counts with
    thousands separators.
  \item \textbf{Zoom and pan.} Drag to zoom, double-click to reset, and a tool bar reduced to
    zoom, pan, reset and download. A redraw with new data keeps the zoom.
  \item \textbf{Colour and shape.} The prototype's blue and orange pair, read from the theme's
    variables so light and dark both work. Unusual is orange and a triangle; not scored is hollow
    grey. Nothing depends on colour alone.
  \item \textbf{Table view.} A switch shows the chart's data as a table, which is also what a
    screen reader gets, with a one-sentence summary above it.
  \item \textbf{Export.} ``Download PNG'' on every chart, named from the chart, the address and the
    snapshot date; the table view offers CSV.
  \item \textbf{States.} A skeleton of the chart's size while loading (no layout jump); an empty
    state with a next action; an error state with ``Try again'' and the table view as the fallback
    if the chart itself fails to draw.
  \item \textbf{Size.} The chart fills its container and redraws when the container resizes; on a
    phone it uses fewer ticks and moves the tools into a menu.
\end{itemize}

\subsection{Each chart}

\subsubsection*{G1 \quad Daily activity}
\begin{itemize}
  \item \textbf{Marks.} One bar per day: transactions that day. A day with an unusual transaction
    is orange with a triangle above it. Days are continuous on the axis, so quiet stretches show as
    gaps; a history shorter than 21 days is padded, as in the prototype.
  \item \textbf{Interaction.} A range slider under the chart and range buttons (1 month, 3 months,
    6 months, 1 year, all); hover for the day's counts; click or Enter opens the day in V3; arrow
    keys move between active days.
  \item \textbf{Edge cases.} One active day; three transactions in all; 348 active days on a phone,
    where the chart opens on the last six months with the slider showing the rest.
\end{itemize}

\subsubsection*{G2 \quad Transactions over time}
\begin{itemize}
  \item \textbf{Marks.} One point per transaction: time on x, the absolute amount in BTC on a log
    y axis. Colour and shape show the status; the size grows with the score. Received and sent can
    be told apart by a filled or an outlined mark and by the filter.
  \item \textbf{Interaction.} Box zoom on both axes; the filter bar (Section~\ref{sec:filters});
    hover for the facts and the main reason; click opens the transaction. With S-07, box or lasso
    selection filters the table to the selected points and shows ``14 selected: 9 received, 5 sent,
    3 unusual''.
  \item \textbf{Edge cases.} A net amount of zero cannot sit on a log axis: it is drawn on the
    axis floor and labelled 0. Amounts spanning eight orders of magnitude are why the axis is log.
  \item \textbf{Scale.} Above \code{webglMinPoints} the chart switches to Plotly's WebGL scatter;
    it falls back to SVG where WebGL is missing.
\end{itemize}

\subsubsection*{G3 \quad Feature deviation}
\begin{itemize}
  \item \textbf{Marks.} Four horizontal bars, one per feature, from 0 to the feature's robust z.
    A shaded band marks $\pm 2$ (stands out) and dashed lines mark $\pm 5.0$ (the flag line, read
    from E2). A bar is labelled with its role in words.
  \item \textbf{Interaction.} Hover: ``Size: 0.05 BTC, against a usual 0.000525 BTC (z 5.9)''.
    In V4 the bars move as the form changes.
  \item \textbf{Edge cases.} A feature with no comparison shows ``no comparison'' in place of a
    bar. A z beyond $\pm 20$ is drawn at the edge with an arrow, since the models clip there.
\end{itemize}

\subsubsection*{G5 \quad What-if response}
\begin{itemize}
  \item \textbf{Marks.} A line: amount on a log x axis, the deciding method's score on y. A
    horizontal line at its threshold; the stretches where it flags are shaded; a vertical line at
    the address's usual amount; a marker at the amount in the form.
  \item \textbf{Interaction.} Dragging the form's slider moves the marker at once; the line
    redraws when the direction or time changes. Hover reads the score at any amount. Clicking the
    curve sets the form's amount. A menu switches the curve to another method, with its own
    threshold line.
  \item \textbf{Why it matters.} The curve shows that the baseline flags unusually small amounts
    as well as unusually large ones, which a single verdict cannot show.
\end{itemize}

\subsubsection*{G4 \quad Model verdicts (S-02)}
\begin{itemize}
  \item \textbf{Marks.} One row per method, grouped by family. The x axis is the share of this
    address's scored transactions with a lower score under that method, from 0 to 100\%. A tick
    shows where the method's threshold falls on the same scale; the dot is an orange triangle when
    the method flags.
  \item \textbf{Why ranks.} The nine scores have different units and one threshold is 0, so
    neither raw scores nor ratios can share an axis. A rank within the address's own history can,
    and it matches what the Explorer means by unusual. E7 and E8 add the ranks when S-02 is built.
\end{itemize}

\subsubsection*{G6 \quad Planted-test recall (S-03)}
\begin{itemize}
  \item \textbf{Marks.} Methods on the x axis; three bars each for the planted patterns (spike,
    burst, dormant) and a marker for the mean with its 95\% interval.
  \item \textbf{Caption.} ``Share of planted test transactions each method flagged in Assignment
    2's evaluation. This is not accuracy on real activity, which has no labels.''
  \item \textbf{Interaction.} Legend entries switch patterns on and off; hover shows the interval.
\end{itemize}

\subsection{Which chart has which interaction}

\begin{center}
\small\setlength{\tabcolsep}{4.5pt}
\begin{tabular}{@{}lcccccc@{}}
\toprule
 & \textbf{G1} & \textbf{G2} & \textbf{G3} & \textbf{G5} & \textbf{G4} & \textbf{G6} \\
\midrule
Zoom and pan & \yes & \yes & -- & \yes & -- & \yes \\
Range slider or range buttons & \yes & -- & -- & -- & -- & -- \\
Tooltips with units & \yes & \yes & \yes & \yes & \yes & \yes \\
Click-through to a view & \yes & \yes & -- & sets the form & -- & -- \\
Filtering & method & filter bar & -- & method & -- & legend \\
Selection (S-07) & -- & \yes & -- & -- & -- & -- \\
Updates without a reload & method, replay & filters, replay & what-if & what-if & what-if & -- \\
Table view, PNG export & \yes & \yes & \yes & \yes & \yes & \yes \\
\bottomrule
\end{tabular}
\end{center}

\subsection{Performance}

\begin{itemize}
  \item \textbf{Loading.} Plotly is the largest dependency. It is imported only when the first
    chart mounts, from a bundle that registers only the trace types in use (bar, scatter, WebGL
    scatter, and heat map if G7 or G8 is built). The first screen never waits for it.
  \item \textbf{Drawing.} Charts update in place and keep the zoom; resizes are debounced; a
    chart off screen is not redrawn.
  \item \textbf{Data.} G2 loads pages of 200 rows and draws when the first page arrives.
  \item \textbf{Evidence (S-12).} \code{tools/measure.py} and a browser script record the time to
    first chart and the redraw time for 700, 7,966 and 50,000 points (the last is synthetic and
    appears in no view). The table goes into the report.
\end{itemize}

\begin{done}
Each finished chart passes the contract's checklist, has a component test for its states and its
table view, and is shown in the video with one zoom, one tooltip and one click. Three chart types
are present at MS2.
\end{done}
